\subsection{AUTOMORPHISMS OF A FRAMED SEMI-SIMPLE LIE ALGEBRA}

Recall (Chap. VII, §3, no. 1) that \(\operatorname{Aut}(\mathfrak{g})\) denotes the group of automorphisms of \(\mathfrak{g}\) . If \(\mathfrak{h}\) is a Cartan subalgebra of \(\mathfrak{g}\) , we denote by \(\operatorname{Aut}(\mathfrak{g},\mathfrak{h})\) the group of automorphisms of \(\mathfrak{g}\) that leave \(\mathfrak{h}\) stable. Assume that \(\mathfrak{h}\) is splitting, and let R be the root system of \((\mathfrak{g},\mathfrak{h})\) . If \(s\in\operatorname{Aut}(\mathfrak{g},\mathfrak{h})\) , the contragredient map of \(s|_{\mathfrak{h}}\) is an element of A(R) (the group of automorphisms of R) which we shall denote by \(\varepsilon(s)\) in this paragraph. Thus

\[
\varepsilon : \operatorname{Aut} (\mathfrak {g}, \mathfrak {h}) \to \mathrm{A} (\mathrm{R})
\]

is a homomorphism of groups.

For any root system R and any basis B of R, we denote by \(\operatorname{Aut}(R,B)\) the group of automorphisms of R that leave B stable. Recall (Chap. VI, §1, no. 5, Prop. 16 and §4, no. 2, Cor. of Prop. 1) that A(R) is the semi-direct product of \(\operatorname{Aut}(R,B)\) and W(R), and that A(R)/W(R) is canonically isomorphic to the group of automorphisms of the Dynkin graph of R.

PROPOSITION 1. Let \((\mathfrak{g},\mathfrak{h},\mathrm{B},(X_{\alpha})_{\alpha \in \mathrm{B}})\) be a framed semi-simple Lie algebra, and R the root system of \((\mathfrak{g},\mathfrak{h})\) . Let G be the set of \(s\in \operatorname {Aut}(\mathfrak{g},\mathfrak{h})\) that leave B stable, and such that \(s(X_{\alpha}) = X_{\varepsilon (s)\alpha}\) for all \(\alpha \in \mathrm{B}\) (in other words the set of automorphisms of \((\mathfrak{g},\mathfrak{h},\mathrm{B},(X_{\alpha})_{\alpha \in \mathrm{B}})\) . Then the restriction of \(\varepsilon\) to \(\mathrm{G}\) is an isomorphism from \(\mathrm{G}\) to \(\operatorname {Aut}(\mathrm{R},\mathrm{B})\) .

If \(s \in G\) , it is clear that \(\varepsilon(s) \in \operatorname{Aut}(\mathrm{R}, \mathrm{B})\) . On the other hand, the map

\[
\varepsilon|_{\mathrm{G}}: \mathrm{G} \rightarrow \operatorname{Aut} (\mathrm{R}, \mathrm{B})
\]

is bijective by Th. 2 of §4, no. 4.

